\chapter{The Almgren--Chriss Framework}\label{ch:mx:the-almgren-chriss-framework}

A trader must sell a million shares by the close. Selling them all now pays the impact; selling them evenly carries the price risk all day. Between
the two lies a curve of schedules, and choosing a point on it is choosing a risk aversion. This chapter sets the problem as Almgren and Chriss did,
solves it in closed form, draws the frontier, and then runs the schedules in the simulated market to see what the model gets right and what it misses.

\section{Parent orders, child orders and a trajectory}

\begin{definition}[Parent order, child order]\label{def:mx:the-almgren-chriss-framework:parent}
A \emph{parent order}\index{parent order} is an instruction to trade a quantity over a period, given to a trader or an algorithm; the
\emph{child orders}\index{child order} are the orders actually sent to venues to execute it.
\end{definition}

\begin{definition}[Trading trajectory]\label{def:mx:the-almgren-chriss-framework:trajectory}
A \emph{trading trajectory}\index{trading trajectory} is the quantity $x(t)$ still to be traded at each time $t$ of the execution window $[0,T]$, from
$x(0)=X$ to $x(T)=0$; its trading rate is $v(t)=-\dot x(t)$.
\end{definition}

The \omterm{def:mx:the-almgren-chriss-framework:parent}{parent order} is a metaorder (chapter~11) seen from the inside. A schedule chooses the trajectory; child-order placement (chapter~17) and routing
(chapter~18) decide how each slice is traded. Bertsimas and Lo (1998) posed the choice as a dynamic programme minimising the expected cost; Almgren and
Chriss (2001) added the risk.

\section{Costs: impact and risk}

\begin{definition}[Almgren--Chriss model]\label{def:mx:the-almgren-chriss-framework:ac}
In the \emph{Almgren--Chriss model}\index{Almgren--Chriss model} the price moves by a permanent impact $\gamma$ per share traded and a random walk of
volatility $\sigma$, and each trade pays a temporary impact $\eta v$ per share at rate $v$. Selling $X$ along $x(t)$ costs, against the arrival price,
\[
\E[C]=\tfrac12\gamma X^2+\eta\int_0^Tv(t)^2\,dt,\qquad \operatorname{Var}[C]=\sigma^2\int_0^Tx(t)^2\,dt.
\]
\end{definition}

\begin{definition}[Execution risk]\label{def:mx:the-almgren-chriss-framework:risk}
\emph{Execution risk}\index{execution risk} is the variance (or standard deviation) of an execution's cost that comes from the price moving while the
position is still held: $\sigma^2\int x^2\,dt$ here.
\end{definition}

The permanent term does not depend on the schedule. The temporary term is smallest for trading as slowly as possible, evenly over the whole window; the
risk is smallest for trading at once. A risk-averse trader minimises $\E[C]+\lambda\operatorname{Var}[C]$ (mean--variance, One Quant Book 7, chapter~25),
with the risk aversion $\lambda$ in units of one over money.

\section{The optimal trajectory in closed form}

\begin{proposition}[Almgren--Chriss trajectory]\label{prop:mx:the-almgren-chriss-framework:sinh}
The trajectory that minimises $\E[C]+\lambda\operatorname{Var}[C]$ is
\[
x(t)=X\,\frac{\sinh\bigl(\kappa(T-t)\bigr)}{\sinh(\kappa T)},\qquad \kappa=\sqrt{\frac{\lambda\sigma^2}{\eta}}.
\]
\end{proposition}

\begin{proof}
The objective is $\int_0^T(\eta\dot x^2+\lambda\sigma^2x^2)\,dt$ plus a constant. Its Euler--Lagrange equation is $\eta\ddot x=\lambda\sigma^2x$, whose
solutions are combinations of $e^{\pm\kappa t}$; the boundary conditions $x(0)=X$, $x(T)=0$ select the ratio of hyperbolic sines. (This is also the value
function route of One Quant Book 4, chapter~9, for a linear-quadratic problem.)
\end{proof}

\begin{definition}[Urgency]\label{def:mx:the-almgren-chriss-framework:urgency}
The \emph{urgency}\index{urgency} of an execution is $\kappa$, or the dimensionless $\kappa T$: zero for a straight line (time-weighted), large for a
trajectory that sells most of the order early. Its inverse is the execution's natural time scale.
\end{definition}

On $n$ intervals of length $\tau$ the discrete optimum has the same form with $\tilde\kappa$ given by $\cosh(\tilde\kappa\tau)=1+\lambda\sigma^2\tau^2/(2\eta)$,
which \texttt{firm.acexec.discrete} implements with the cost and variance of any discrete schedule.

\omcode{code/firm/acexec/firm_acexec.py}{40}{57}{The discrete Almgren--Chriss trajectory and the expected cost and variance of a schedule.}\label{lst:mx:the-almgren-chriss-framework:ac}

The chapter's block: a stock at USD~50 trading 10 million shares a day with 2\% daily volatility, so $\sigma=1$ dollar per share per square root of a day;
selling 1 million by the close. The temporary impact is linearised from chapter~11's square-root law at 10\% participation,
$\eta=0.8\times0.02\times\sqrt{0.1}\times50/(0.1\times10^7)=2.53\times10^{-7}$ dollars per share per share a day; no permanent impact. With $\lambda=10^{-6}$ per
dollar, $\kappa=1.99$ per day. The optimal schedule on one-minute intervals expects to cost USD~300\,000 (30 cents a share, 60 basis points) with a
standard deviation of USD~470\,000, a certainty equivalent $\E+\lambda\operatorname{Var}$ of USD~521\,000, and sells 95\% of the order in 5.9 of the 6.5 hours
(\cref{fig:mx:the-almgren-chriss-framework:frontier}, left).

\section{The efficient trading frontier}

\begin{definition}[Efficient trading frontier]\label{def:mx:the-almgren-chriss-framework:frontier}
The \emph{efficient trading frontier}\index{efficient trading frontier} is the set of (risk, expected cost) pairs of the optimal trajectories as the risk
aversion varies: no schedule has both lower risk and lower expected cost than a point on it.
\end{definition}

The straight line is the frontier's end with the least expected cost: USD~253\,000 with a standard deviation of USD~576\,000 and a certainty equivalent of
USD~585\,000 at $\lambda=10^{-6}$. What does getting the \omterm{def:mx:the-almgren-chriss-framework:urgency}{urgency} wrong cost? A schedule twice too patient ($\kappa/2$) expects USD~258\,000 with a standard
deviation of 542\,000, a certainty equivalent of 551\,000, 5.8\% worse than the optimum; one twice too hurried ($2\kappa$) expects 506\,000 with 352\,000, a
certainty equivalent of 630\,000, 20.9\% worse. The frontier is flat near its optimum and steep toward hurry: erring on the side of patience is cheaper.

\begin{omfigure}
\begin{tikzpicture}
\begin{groupplot}[group style={group size=2 by 1,horizontal sep=1.6cm},width=6.6cm,height=5cm,
  tick label style={font=\footnotesize},label style={font=\small},grid=major,grid style={gray!25},table/col sep=comma]
\nextgroupplot[title={\small share of the order left},xlabel={hours},xmin=0,xmax=6.5,ymin=0,ymax=1,
  legend columns=2,legend style={font=\footnotesize,at={(0.5,-0.3)},anchor=north,draw=none,fill=none}]
\addplot[black,thick] table[x=hours,y=k0]{figdata/microstructure/14-the-almgren-chriss-framework/trajectories.csv};
\addplot[omProp,thick] table[x=hours,y=k1]{figdata/microstructure/14-the-almgren-chriss-framework/trajectories.csv};
\addplot[omDef,very thick] table[x=hours,y=k2]{figdata/microstructure/14-the-almgren-chriss-framework/trajectories.csv};
\addplot[omThm,thick] table[x=hours,y=k3]{figdata/microstructure/14-the-almgren-chriss-framework/trajectories.csv};
\legend{$\kappa=0$,$\kappa=1$,$\kappa=1.99$,$\kappa=8$}
\nextgroupplot[title={\small efficient frontier (USD thousands)},xlabel={standard deviation},ylabel={expected cost},xmin=0,xmax=620,ymin=150,ymax=700,
  legend columns=2,legend style={font=\footnotesize,at={(0.5,-0.3)},anchor=north,draw=none,fill=none}]
\addplot[omDef,very thick] table[x=sd,y=cost]{figdata/microstructure/14-the-almgren-chriss-framework/frontier.csv};
\addplot[omThm,only marks,mark=*] table[x=sd,y=cost]{figdata/microstructure/14-the-almgren-chriss-framework/points.csv};
\node[omThm,font=\scriptsize,anchor=west] at (axis cs:360,506) {hurried};
\node[omThm,font=\scriptsize,anchor=north east] at (axis cs:462,292) {optimal};
\node[omThm,font=\scriptsize,anchor=north] at (axis cs:535,240) {patient};
\node[omThm,font=\scriptsize,anchor=south] at (axis cs:585,262) {line};
\end{groupplot}
\end{tikzpicture}

\omcaption{Selling 1 million shares by the close. Left: optimal trajectories for several urgencies (per day). Right: the efficient frontier, with the
optimum for $\lambda=10^{-6}$, the time-weighted schedule, and schedules twice too patient and twice too hurried. Data: \texttt{mx\_ac.block\_study},
\texttt{fig\_ac.py}.}\label{fig:mx:the-almgren-chriss-framework:frontier}
\end{omfigure}

\section{Calibration and limits}

The model has three inputs, and each is an estimate. The execution agent of \texttt{mx\_ac} runs a scheduler in \texttt{firm.agentmkt}'s market: it sells
20\,000 shares over 30 minutes in 30 slices, each down to the scheduler's target, for urgencies $\kappa T$ of 0, 2 and 6, on eight sessions each. The
scheduler interface of \texttt{firm.acexec} (target holdings at given times) is the one chapters 15, 16 and 28 implement.

\omcode{code/microstructure/14-the-almgren-chriss-framework/python/mx_ac.py}{55}{71}{An execution agent: at each interval it sells down to the scheduler's target.}\label{lst:mx:the-almgren-chriss-framework:agent}

The realised shortfall per share (\cref{fig:mx:the-almgren-chriss-framework:sim}) is 5.2 ticks for the straight line (standard error 2.8), 5.3 (2.3) for
$\kappa T=2$ and 7.4 (1.3) for $\kappa T=6$, with standard deviations across sessions of 8.0, 6.5 and 3.8: faster schedules cost more on average and vary
less, as the model says. Calibrated on the same runs, the temporary impact is $\eta=0.11$ ticks per share per share a second (from each slice's cost
against its size) and $\sigma=0.69$ ticks per root second (from one-minute mid changes), and the model predicts means of 1.2, 1.4 and 3.6 ticks and standard
deviations of 16.4, 13.2 and 7.6. It gets the ordering right and the levels wrong in two ways. The mean is too low because each slice's slippage misses the
impact that builds up across slices: the gap implies a permanent impact of about $4\times10^{-4}$ ticks per share, the $\gamma$ set to zero here. The
standard deviation is too high because the market mean-reverts: its fundamentalists pull the price back, so thirty minutes' variance is less than thirty
times one minute's. Almgren (2003) extended the model to non-linear impact and to risk that grows with trading; chapter~15 lets the schedule react to the
market.

\begin{omfigure}
\begin{tikzpicture}
\begin{groupplot}[group style={group size=2 by 1,horizontal sep=1.6cm},width=6.6cm,height=5cm,
  tick label style={font=\footnotesize},label style={font=\small},grid=major,grid style={gray!25},table/col sep=comma,
  xlabel={urgency $\kappa T$},xtick={0,2,6},xmin=-0.5,xmax=6.5]
\nextgroupplot[title={\small mean shortfall (ticks a share)},ymin=0,ymax=11,
  legend columns=1,legend cell align=left,legend style={font=\footnotesize,at={(0.5,-0.3)},anchor=north,draw=none,fill=none}]
\addplot[omDef,very thick,mark=*,error bars/.cd,y dir=both,y explicit] table[x=kt,y=mean,y error=se]{figdata/microstructure/14-the-almgren-chriss-framework/sim.csv};
\addplot[omThm,dashed,mark=square*] table[x=kt,y=model_mean]{figdata/microstructure/14-the-almgren-chriss-framework/sim.csv};
\legend{simulated,model}
\nextgroupplot[title={\small standard deviation (ticks a share)},ymin=0,ymax=18]
\addplot[omDef,very thick,mark=*] table[x=kt,y=sd]{figdata/microstructure/14-the-almgren-chriss-framework/sim.csv};
\addplot[omThm,dashed,mark=square*] table[x=kt,y=model_sd]{figdata/microstructure/14-the-almgren-chriss-framework/sim.csv};
\end{groupplot}
\end{tikzpicture}

\omcaption{Almgren--Chriss schedules executed in the simulated market (eight sessions each) against the model calibrated on the same runs. The model
orders the schedules correctly; it misses the impact that accumulates across slices and the market's mean reversion. Data: \texttt{mx\_ac.sim\_study}.}\label{fig:mx:the-almgren-chriss-framework:sim}
\end{omfigure}

\section{Tutorial: liquidating a block by the close}\label{tut:mx:the-almgren-chriss-framework:tut}

\begin{tutorial}
\textbf{Goal.} Solve the Almgren--Chriss problem for a block, draw its frontier, and test the schedules in the simulated market.
\textbf{End state:} \cref{fig:mx:the-almgren-chriss-framework:frontier,fig:mx:the-almgren-chriss-framework:sim}.
\begin{tutsteps}
\item \textbf{Model.} \texttt{firm\_acexec.kappa}, \texttt{trajectory}, \texttt{discrete}, \texttt{cost\_var}, \texttt{frontier}, \texttt{time\_to}.
\item \textbf{Block.} \texttt{mx\_ac.block\_study(lam)}: the optimum, the time-weighted schedule and the two mistakes.
\item \textbf{Simulation.} \texttt{ScheduleAgent(ACScheduler(Q, 1, kT), \dots)} in \texttt{firm\_agentmkt.session}; \texttt{sim\_study()} for three
urgencies and eight seeds, with the calibration.
\item \textbf{Draw.} \texttt{fig\_ac.py}.
\end{tutsteps}
\textbf{What to change next.} Add the permanent impact implied by the simulation to the model and recompute its means; calibrate $\sigma$ at the
execution's horizon instead of one minute.
\end{tutorial}

\section{Build: the execution schedule}\label{bld:mx:the-almgren-chriss-framework:acexec}

\begin{build}
\bfield{Purpose} The reference schedule of the firm's execution: the benchmark that chapters 15 and 16 improve on, the scheduler interface their
algorithms and chapter~28's implement, and the pre-trade cost and risk estimates of chapter~19.
\bfield{Interface} \texttt{kappa(eta, sigma, lam)}, \texttt{trajectory(X, T, kappa, t)}, \texttt{discrete(X, T, n, eta, sigma, lam)},
\texttt{cost\_var(holdings, T, gamma, eta, sigma)}, \texttt{frontier(\dots, lams)}, \texttt{time\_to(share, kappa, T)}; \texttt{Scheduler.targets(times)},
\texttt{ACScheduler}, \texttt{LinearScheduler}.
\bfield{Rules} Holdings are the quantity still to trade; the discrete \omterm{def:mx:the-almgren-chriss-framework:urgency}{urgency} from $\cosh(\tilde\kappa\tau)$; costs against the arrival price; $\eta$
per unit of trading rate in the caller's time unit.
\bfield{Acceptance tests} \texttt{code/firm/acexec/tests/}: discrete and continuous trajectories agree on a fine grid; cost and variance of a straight line
by hand; the frontier's costs rise and variances fall with $\lambda$; the time to 95\% equals $\ln 20/\kappa$ on a long horizon; the schedulers' targets.
\bfield{Stretch} Non-linear temporary impact (Almgren, 2003); a portfolio version with a covariance matrix; calibration from \texttt{firm.impactfit}.
\end{build}

\begin{omsources}
\item D.~Bertsimas and A.~W.~Lo, ``Optimal control of execution costs'', \emph{Journal of Financial Markets} 1(1), 1998.
\item R.~Almgren and N.~Chriss, ``Optimal execution of portfolio transactions'', \emph{Journal of Risk} 3(2), 2001.
\item R.~Almgren, ``Optimal execution with nonlinear impact functions and trading-enhanced risk'', \emph{Applied Mathematical Finance} 10(1), 2003.
\end{omsources}

\section{Exercises}

\begin{exercise}[$\star$]\label{exo:mx:the-almgren-chriss-framework:1}
With $\eta=2.53\times10^{-7}$ dollars per share per share a day, $\sigma=1$ dollar per share per root day and $\lambda=10^{-6}$ per dollar, compute $\kappa$.
\end{exercise}

\begin{exercise}[$\star$]\label{exo:mx:the-almgren-chriss-framework:2}
Compute the expected cost and standard deviation of selling 1 million shares in a straight line over one day with these parameters and no permanent impact.
\end{exercise}

\begin{exercise}[$\star$]\label{exo:mx:the-almgren-chriss-framework:3}
How long does the optimal schedule take to sell 95\% of the order, and what is the formula when the window is long?
\end{exercise}

\begin{exercise}[$\star\star$]\label{exo:mx:the-almgren-chriss-framework:4}
Why does a schedule twice too hurried lose more certainty equivalent (20.9\%) than one twice too patient (5.8\%)?
\end{exercise}

\begin{exercise}[$\star\star$]\label{exo:mx:the-almgren-chriss-framework:5}
Why does the model calibrated from slice slippage predict too low a mean shortfall in the simulation?
\end{exercise}

\begin{exercise}[$\star\star$]\label{exo:mx:the-almgren-chriss-framework:6}
Why does it predict too high a standard deviation?
\end{exercise}

\begin{exercise}[$\star\star\star$]\label{exo:mx:the-almgren-chriss-framework:7}
\emph{Coding.} From the simulated and model means, infer the permanent impact $\gamma$ that would close the gap for $\kappa T=0$ and for $\kappa T=6$. Are the
two estimates consistent?
\end{exercise}

\begin{exercise}[$\star\star\star$]\label{exo:mx:the-almgren-chriss-framework:8}
\emph{Find the flaw.} ``A time-weighted schedule has no risk because it trades the same amount every minute.''
\end{exercise}

\section{Problem: Liquidating a Block by the Close}

\begin{problem}[{Weekend problem --- liquidating a block by the close}]\label{pb:mx:the-almgren-chriss-framework:1}
A portfolio manager must sell 1 million shares of a 50-dollar stock by the close and asks the desk for a schedule, its expected cost and its risk.

\medskip\noindent\textbf{Part I --- The model.}
\begin{enumerate}
\item Define parent and \omterm{def:mx:the-almgren-chriss-framework:parent}{child orders} and a \omterm{def:mx:the-almgren-chriss-framework:trajectory}{trading trajectory}.
\item Write the Almgren--Chriss expected cost and variance.
\item Why does the permanent term not depend on the schedule?
\item Derive the optimal trajectory.
\end{enumerate}

\medskip\noindent\textbf{Part II --- The block.}
\begin{enumerate}[resume]
\item How is $\eta$ linearised from the square-root law, and what is it here?
\item Give $\kappa$, the expected cost, the standard deviation and the certainty equivalent at $\lambda=10^{-6}$.
\item When is 95\% of the order sold?
\item Give the straight line's cost, risk and certainty equivalent.
\end{enumerate}

\medskip\noindent\textbf{Part III --- The frontier.}
\begin{enumerate}[resume]
\item Define the \omterm{def:mx:the-almgren-chriss-framework:frontier}{efficient trading frontier}.
\item Give the cost, risk and certainty equivalent of the schedules twice too patient and twice too hurried.
\item Which mistake is cheaper, and why?
\item What does the risk aversion mean for the manager?
\end{enumerate}

\medskip\noindent\textbf{Part IV --- Calibration and the verdict.}
\begin{enumerate}[resume]
\item Describe the simulated test.
\item Give the realised means and standard deviations.
\item Give the model's calibrated parameters and predictions.
\item Explain the gaps.
\item What permanent impact closes the mean gap?
\item State the \emph{named result}: the optimal trading horizon, expected cost and standard deviation for a given risk aversion, and the extra cost of a
schedule that is twice too patient or twice too hurried.
\item What would you tell the manager about the numbers' reliability?
\item In one sentence: what does the frontier add to a cost estimate?
\end{enumerate}
\end{problem}

\section{Interview questions}

\begin{interviewq}[$\star$ \iqroles{trader, researcher}]\label{iq:mx:the-almgren-chriss-framework:1}
Explain the trade-off in optimal execution and the shape of the Almgren--Chriss trajectory.
\end{interviewq}

\begin{interviewq}[$\star\star$ \iqroles{researcher}]\label{iq:mx:the-almgren-chriss-framework:2}
Derive the Almgren--Chriss trajectory from its objective.
\end{interviewq}

\begin{interviewq}[$\star\star$ \iqroles{trader}]\label{iq:mx:the-almgren-chriss-framework:3}
Your client doubles its risk aversion. What happens to the schedule, the expected cost and the risk?
\end{interviewq}

\begin{interviewq}[$\star\star$ \iqroles{researcher}]\label{iq:mx:the-almgren-chriss-framework:4}
How would you calibrate the temporary and permanent impact of the model?
\end{interviewq}

\begin{interviewq}[$\star\star$ \iqroles{developer}]\label{iq:mx:the-almgren-chriss-framework:5}
Design a scheduler interface that different execution algorithms can share. What does it take and return?
\end{interviewq}

\begin{interviewq}[$\star\star\star$ \iqroles{researcher}]\label{iq:mx:the-almgren-chriss-framework:6}
What does the \omterm{def:mx:the-almgren-chriss-framework:ac}{Almgren--Chriss model} leave out, and which of these matters most for a liquid stock?
\end{interviewq}
